\documentclass[10pt]{article}
\usepackage[margin=20mm]{geometry}
\usepackage{amsmath,amssymb,mathtools,fontspec,unicode-math,xurl,hyperref}
\IfFontExistsTF{DejaVu Serif}{\setmainfont{DejaVu Serif}}{\setmainfont{Latin Modern Roman}}
\setmathfont{latinmodern-math.otf}
\setlength{\parindent}{0pt}\setlength{\parskip}{6pt}\sloppy
\title{Technical review of
Singular light leaves}
\author{AI-assisted technical review}
\date{}
\begin{document}
\maketitle
Abstract: This review concerns the acquired 67-page arXiv:2401.03053v1 manuscript. One substantive Minor finding is established: the last equality of Lemma 8.7 confuses graded rank with graded dimension. An admissible rank-one example gives different degree-two coefficients. The interchanged left and right redundancy labels in Definition 4.15 are retained as a scope clarification, without a demonstrated affected calculation. The corrected forms preserve the subsequent bimodule construction and counting argument. Independent finite-group and Frobenius-algebra calculations support selected interfaces of the proof. No Major defect has been established in the basis theorems. This is a bounded initial review. The retained source evidence images have now been inspected; several external results remain dependencies, and scientific adjudication remains pending.\par
\section*{Source identity and precise scope}
The source is B. Elias, H. Ko, N. Libedinsky and L. Patimo, Singular light leaves, arXiv:2401.03053v1 [source]. The arXiv submission is 5 January 2024 at 20:09:40 UTC; the acquired PDF’s first page separately says 9 January 2024. Both dates are retained. The associated journal DOI is 10.4171/JEMS/1715, with online publication on 24 October 2025 according to EMS Press [ems]. The acquired preprint has not been verified as the accepted manuscript or journal version. All findings here apply only to its exact bytes. The frozen identifier M25-0338 and assigned year 2025 are unchanged.\par
All 67 PDF pages were read through the retained Poppler text extraction. Page references below are one-based PDF pages, agreeing with printed pages. The review followed the double-coset reduction, Hecke coefficients, grading conventions, Frobenius maps, construction of sprinkles, dual sequences, and the independence–spanning argument. Many formulas in Sections 5–8 are diagrams. The bounded continuation visually inspected all 53 retained source evidence pages, including the relevant oriented crossings and proof diagrams. This confirms the displayed notation and cited interfaces; it does not independently reprove the full Frobenius cube calculus.\par
Theorem 7.49, p. 47, asserts a basis after applying the evaluation functor to singular Bott–Samelson bimodules. It is not a complete generators-and-relations presentation of the diagrammatic category. The coefficient rings \(R_p\) vary with the double coset; the displayed global basis is a \(\Bbbk\)-basis after choosing homogeneous bases of these rings. These restrictions are essential to a faithful assessment.\par
\section*{Hypotheses, normalization, and proof structure}
Section 4.1, p. 17, fixes a field and a faithful, balanced, reflection-faithful realization with generalized Demazure surjectivity. All parabolics in the expressions are finitary. The manuscript explicitly warns that the required realization need not exist in positive characteristic for an infinite Coxeter group. Nothing in this review extends its theorem to those excluded realizations.\par
Write \(\underline p,\overline p\) for minimal and maximal representatives, and \(\ell(I)=\ell(w_I)\). These decorations must be kept distinct: text extraction sometimes collapses them. The source’s conventions are \[R=\operatorname{Sym}(V),\qquad \deg V=2,\qquad
 M(n)_j=M_{j+n}.\] For a multistep expression with intermediate maxima \(K_j\) and minima \(I_j\), its Bott–Samelson shift is \[\sigma=\sum_j\ell(K_j)-\sum_{j\geq1}\ell(I_j).\] Thus an increasing step \(I\subset K\) gives \(R^I\) without a shift, whereas the reverse step gives \(R^I(\ell(K)-\ell(I))\). The duality is \(\operatorname{Hom}_{R^I}(-,R^I)\). This agrees with the comprehensive normalization in Williamson’s appended erratum, arXiv:1010.1283v3, Section 8.4 [williamson]. The erratum does not by itself invalidate this manuscript’s normalization.\par
For \(p\subset q\), where the right parabolic changes from \(J\) to \(Js\), the upward defect is \[d_\uparrow=\ell(\underline q)-\ell(\underline p)
             -\ell(\operatorname{LR}(q))+\ell(\operatorname{LR}(p)),
 \qquad
 d_\downarrow=d_\uparrow+\ell(Js)-\ell(J).\] The upward polynomial is \(\pi^+(\operatorname{LR}(q))/\pi^+(\operatorname{LR}(p))\), and the downward polynomial is \(1\). Grassmannian reduction isolates the nontrivial local pair and preserves these data. Summing over subordinate paths gives the standard-basis coefficient \(c_p(v)\). With the target self-dual, the Hom formula uses \[\operatorname{gd}\operatorname{Hom}(B,B')=\sum_p c_p(v)c'_p(v)\,\operatorname{gd}R_p .\] In particular, the polynomial factors must not be dropped.\par
The local Frobenius pairing evaluates a light leaf with a dual sprinkle to a Kronecker delta modulo positive-degree elements. With homogeneous choices \(P_L\) of degree \(2\ell(w_L)\), degree comparison removes the possible remainder in the diagonal and triangular cases. Such choices exist: take that homogeneous component of any preimage of \(1\) under the degree \(-2\ell(w_L)\) Demazure operator. This is the graded choice used in the audit, not a verification for arbitrary inhomogeneous replacements.\par
The dual sequence is injective on coterminal sprinkled paths. At the last place where their underlying cosets differ, that step must be increasing; equality of the dual elements would put the two Grassmannian minima in the same double coset, a contradiction. Lexicographic comparison then yields Proposition 8.26. Evaluating a nonzero combination on the maximal dual test tensor detects its coefficient. Lower-term separation, followed by the corrected graded count, gives the basis modulo lower terms. Finally, induction on the finite Bruhat support gives spanning of double leaves; the full Hom dimension gives independence. The two-sided triangular order, rather than either individual total order alone, controls Theorem 7.51.\par
\section*{Scope clarification: the redundancy inclusions}
Definition 4.15, p. 21, writes \(K=\operatorname{LR}(p)\subset J\) and \(L=\operatorname{RR}(p)\subset I\). These containing labels are interchanged. The following example confirms the textual correction; it is not counted as a substantive Problem. Take type \(A_2\) over \(\mathbb Q\) with its usual reflection realization, \(s=(12)\), \(t=(23)\), \(I=\{s\}\), \(J=\{t\}\), and \(w=ts\). This realization satisfies the stated hypotheses, and all the parabolics are finite. With composition acting rightmost first, \[W_IwW_J=\{ts,sts\},\qquad \ell(ts)=2,\quad\ell(sts)=3.\] Hence \(w=\underline p\). Direct conjugation gives \[wtw^{-1}=s,\qquad w^{-1}sw=t,
 \qquad
 \operatorname{LR}(p)=I,\quad \operatorname{RR}(p)=J.\] Since \(s\ne t\), neither printed cross-inclusion holds. The retained permutation calculation reproduces the coset and its two lengths exactly.\par
The correct statements are \[K=\operatorname{LR}(p)\subset I,\qquad L=\operatorname{RR}(p)\subset J.\] They follow directly from Section 2.2, p. 10. They also give exactly the ring inclusions needed by the existing construction: \[R^I\subset R^K,\qquad
 R^J\subset R^L,\qquad
 \underline p R^L=R^K.\] Thus left multiplication by \(R^I\) and the right action \(m\cdot f=(\underline p f)m\) preserve \(R^K\) and commute. In the displayed example \(R_p=R^I\), and every left-linear endomorphism is multiplication by its value at \(1\); it automatically commutes with the twisted right action. Consequently \(\operatorname{End}(R_p)=R_p\) here.\par
Impact and repair boundary. This is a scope clarification of the containing labels. The named redundancy subgroups are correctly defined on p. 10 and are the subgroups actually used in the subsequent construction. No changed action, calculation, realization assumption, basis element, or conclusion has been demonstrated. The annotation error alone is therefore excluded from the substantive finding count. The corrected labels agree with the existing actions; no failure of the standard-bimodule theory is inferred.\par
\section*{Graded rank versus graded dimension}
Problem 1 (A missing polynomial-ring factor in (179)).  The last equality of Lemma 8.7, p. 52, equates the finite light-leaf degree polynomial with \(\operatorname{gd}\operatorname{Hom}_{\not<p}(\operatorname{BS}(I_\bullet),\operatorname{BS}(Y_p))\). With the full graded-dimension convention used in Theorem 4.16, this equality is false.\par
Use type \(A_1\) over \(\mathbb Q\), \(R=\mathbb Q[x]\), \(\deg x=2\), and the identity expression \(I_\bullet=Y_p=[\varnothing]\), \(p=e\). There is one subordinate path, no sprinkles, and one leaf, the identity of degree zero. There is no lower Bruhat layer. Therefore \[\sum_{\mathrm{LL}}v^{\deg\mathrm{LL}}=1,\qquad
 \operatorname{Hom}_{\not<e}(R,R)=R,\qquad
 \operatorname{gd}R=\frac{1}{1-v^2}.\] The degree-two coefficients are respectively \(0\) and \(1\). This is an exact counterexample under all the paper’s hypotheses; it is not a finite numerical extrapolation or a grading-shift ambiguity.\par
Proposition 4.21, p. 22, already supplies the repair: \(\operatorname{Hom}_{\not<p}\cong c_p(v)R_p\). In a notation which removes all ambiguity, (179) should conclude \[\sum_{\mathrm{LL}}v^{\deg\mathrm{LL}}
 =\sum_{t_\bullet}\operatorname{poly}(t_\bullet)
                    v^{\operatorname{def}(t_\bullet)}
 =\frac{\operatorname{gd}\operatorname{Hom}_{\not<p}(\operatorname{BS}(I_\bullet),\operatorname{BS}(Y_p))}
        {\operatorname{gd}R_p}
 =c_p(v).\] The quotient is the intended graded rank for the indicated free-module identification. Equivalently, multiply the two sums by \(\operatorname{gd}R_p\) when writing a full graded-dimension equality.\par
Actual downstream effect. Lemma 8.8, p. 53, explicitly restores the factor \(\sum_{b\in B_p}v^{\deg b}=\operatorname{gd}R_p\). Proposition 8.32, p. 63, expressly compares graded ranks. Thus the main counting argument uses the corrected quantity. This is assessed as Minor because the proof’s needed count is recovered under the same hypotheses and its later use is consistent with that count. It is not evidence that the full polynomial multiples of leaves are absent from the proposed basis.\par
\section*{Independent checks and a proof interface}
\subsection*{Finite Coxeter calculations}
The retained independent Python program enumerates all 555 pairs relevant to Grassmannian reduction in \(S_4\), over every pair of standard parabolics and every allowed added simple reflection. It verifies the minimal-element factorization \(\underline p=\underline q\,\underline m\), length additivity, redundancy conjugacy, reduced coset lengths, and both defect identities. A separate Hecke multiplication implementation in \(S_3\), using \[h_s^2=1+(v^{-1}-v)h_s,\] checks 42 upward and 24 downward instances of Proposition 3.4. It also reproduces the six distinct dual sequences in Example 8.21 and the order in Example 8.24: \[sts\prec ts\prec s\prec st\prec t\prec e .\] All assertions passed. These finite checks support the stated interfaces; they do not prove the arbitrary-Coxeter theorem.\par
\subsection*{Rank-one Frobenius and Hom calculations}
Let \(s(x)=-x\), \(A=\mathbb Q[x^2]\), and \(\partial f=(f-sf)/(2x)\), so \(\partial x=1\). For \(B=R\otimes_A R(1)\), use the left \(R\)-basis \(e=1\otimes1,\ f=1\otimes x\). Right multiplication by \(x\) has matrix \[X=\begin{pmatrix}0&x^2\\1&0\end{pmatrix}.\] A left-linear matrix commutes with \(X\) precisely when it is \[M=\begin{pmatrix}a&x^2c\\c&a\end{pmatrix}.\] Writing \(\epsilon\) for multiplication and \(\eta(1)=x\otimes1+1\otimes x\), we have \[T=\eta\epsilon=x\,\mathrm{id}+X,\qquad
 T^2=2xT,\qquad
 \operatorname{End}_{R-R}(B)=R\,\mathrm{id}\oplus R\,T.\] The basis degrees are \(0,2\), and the graded rank is \(1+v^2\).\par
For the singular object \(R(1)\) over \((A,A)\), the correct answer has a negative-degree term: \[\operatorname{End}_{A-A}(R(1))
 =R\,\mathrm{id}\oplus R\,\partial,\qquad
 \operatorname{gd}\operatorname{End}_{A-A}(R(1))
 =(1+v^{-2})\operatorname{gd}R=(v+v^{-1})^2\operatorname{gd}A.\] Indeed, a map taking \(1\) to \(u\) and \(x\) to \(z\) is \(u\,\mathrm{id}+(z-xu)\partial\). The four \(A\)-matrix units have degrees \(0,0,2,-2\). A common shift of source and target does not change them. These formulas agree with the manuscript’s shifts and Hom formula. An initial analytical test expectation for this singular case was corrected during the bounded auxiliary discussion before code execution; it was not treated as a paper defect.\par
\subsection*{The lower-term separation used in Theorem 8.31}
Proposition 4.20 is stated with a reduced source as well as a reduced target, whereas Theorem 8.31 applies the separation property to arbitrary source expressions. The needed extension can be checked without assuming the light-leaves theorem. The cited companion paper’s Theorem 3.30 identifies lower maps with maps into \(\Gamma_{<p}\) of the target. Its Proposition 3.34 proof constructs the target quotient \[q_p:\operatorname{BS}(Y_p)\longrightarrow R_p(\text{grading shift}),
 \qquad q_p(h1^\otimes)=h.\] Thus, for any source \(M\) and lower map \(F:M\to\operatorname{BS}(Y_p)\), \(q_pF=0\). If \(F(z)=h1^\otimes\), then \(h=0\). The source need not be reduced. The companion’s exact sequences and this implication were checked in arXiv:2307.15726v2 [lower]. This supplies the required interface; no counterexample to the independence argument follows from the abbreviated citation.\par
\section*{Coverage, dependencies, and assessment}
The textual chain through Sections 2–8 has been followed, including the well-founded Bruhat induction, the distinction between a reduced target and an arbitrary expression, and the final cellular packaging. The local dual-basis hypotheses in Theorem 8.15 were compared with Lemma 4.13 and Theorem 4.14 of arXiv:2307.15021v2 [demazure]. The relevant lower-term quotient statements and the normalization in Williamson’s erratum were checked at their cited primary sources. Their complete foundational proofs, the full Frobenius cube calculus, and every diagram in this 67-page manuscript have not been independently reproved. The cited source evidence diagrams have now been visually inspected; external primary-source consultations are retained facts of the original execution and are not credited as new image inspections.\par
For a homogeneous free-module map with the stated independent generators and the same graded ranks, equality of the finite-dimensional homogeneous pieces supplies spanning. In the last Bruhat induction one may use a finite downward-closed support, as Remark 8.36 explicitly allows; no single maximal element of an arbitrary infinite Coxeter group is needed. The cellular statement concerns the category of Bott–Samelson objects in the theorem’s setting and uses the chosen identity leaf for a reduced object. It is not asserted here for an unspecified enlarged category.\par
The bounded source search checked the arXiv record, publisher metadata, an author publication page, and correction/dependency records. No separate paper-specific code release or supplement was located. The publicly linked companion papers are mathematical dependencies. The publisher page requires a subscription for the article, so the journal text was not acquired through this review. Initial direct versioned dependency requests returned tool errors; their public author links led to the accessible v2 HTML records used above. The initial publisher-erratum fetch failed; the author’s link to the appended arXiv erratum supplied the normalization check. These limitations are retained, not interpreted as absence of a correction. No author code was run.\par
The overall finding is Minor, for the graded-count discrepancy in Lemma 8.7. The redundancy-label interchange is a separate scope clarification, not a substantive finding. No Major issue has been established within the actual coverage. That statement is not an independent acceptance or an unconditional certification of the basis theorem. The source evidence images have been inspected in the bounded continuation; the exact final report build and page inspections are recorded separately in the recovery evidence.\par
9 B. Elias, H. Ko, N. Libedinsky and L. Patimo, Singular light leaves, arXiv:2401.03053v1, 5 January 2024; PDF date 9 January 2024. https://arxiv.org/abs/2401.03053v1. EMS Press, article metadata for DOI 10.4171/JEMS/1715, consulted 27 October 2026. https://ems.press/journals/jems/articles/14299228. G. Williamson, Singular Soergel bimodules, arXiv:1010.1283v3, 2 January 2024, appended erratum, Sections 8.4–8.6. https://arxiv.org/html/1010.1283v3. B. Elias, H. Ko, N. Libedinsky and L. Patimo, Subexpressions and the Bruhat order for double cosets, arXiv:2307.15726v2, 5 January 2024, Section 3.8. https://arxiv.org/html/2307.15726v2. B. Elias, H. Ko, N. Libedinsky and L. Patimo, Demazure operators for double cosets, arXiv:2307.15021v2, 5 January 2024, Section 4. https://arxiv.org/html/2307.15021v2.\par
\section*{Supporting checks}
Paper identifier: `M25-0338`.\par
618d2b9da220747cb931ba9e129502a9b68be265b50a3e1a44647276adb77e6b\par
The initial executor read the text, performed the independent calculations, assessed the findings, and wrote this report. The existing native auxiliary agent `/root/bounded\_0331\_check` challenged the supplied \(A_2\) redundancy example, rank-one endomorphism calculations, the rank/dimension test, and the target-quotient argument. It did not read local files, execute code, inspect images, or conduct independent acceptance. The executor’s subsequent dependency reading and computations are separate from that analytical contribution.\par
The retained independent program completed with exit code \(0\). The existing official Tectonic 0.17.0 route is used with this paper’s own cache and temporary directory. Neither rendering nor parsing is credited as viewing. Earlier EMFILE image-access failures and the original attachment-needed record remain preserved as historical evidence. Current source and report inspection receipts separately identify the images actually viewed.\par
The native bounded continuation inspected all 53 retained source evidence pages, confirmed the two local findings and their downstream limits, and retained the original calculation results without rerunning them. Its exact-source final build and all report-page inspections are recorded under this paper’s recovery evidence, with the original execution origins kept distinct. Scientific decision remains `pending`, with no human review or independent acceptance recorded. Runtime self-description and the controller’s reported official UI selection are separate records; no backend-model equivalence or personal UI observation is asserted.\par
\end{document}
